\title{XQDot4Z: Quantized Dot Products\\on a RISC-V Processor}
\author{\IEEEauthorblockN{Steve Andy Ildefonso Santos}
\IEEEauthorblockA{Computer Science student\\
Working draft v\PaperVersion\ --- not for submission}}
\maketitle
\begin{abstract}
Low-bit-weight storage reduces data size, but processing these weights may
require extraction, conversion, and correction of the numerical origin.
This work proposes studying an operation that combines four products of
unsigned four-bit weights and signed eight-bit activations with zero-point
correction. The question is when this fusion is useful compared with scalar
multiplication and with a packed dot product followed by separate correction.
A teaching-oriented thirty-two-bit processor is proposed as the platform,
first to verify the operation and then to evaluate dot products and
matrix--vector multiplication. This document is an initial draft: it presents
the question and reports a registered campaign of 2280 simulated cases. Fusion
reduces cycles against both alternatives, by a margin that narrows as the row
count grows and that vanishes when the zero point is zero. With a varying zero
point the advantage is paid for in code size, because the immediate forces
specialization. These are simulated cycles: not physical time, energy or area.
\end{abstract}
\begin{IEEEkeywords}
quantized inference, low-precision arithmetic, instruction-set extensions,
hardware verification
\end{IEEEkeywords}
% Keep the abstract first; a top float may start in the next column instead.
\suppressfloats[t]

\section{Introduction}
Neural network layers perform products and sums involving weights and
activations. Quantization allows part of this computation to be represented
using integers and scale parameters~\cite{jacob}. However, storing more weights
per word does not eliminate the work needed to consume them on a processor.
Field extraction and operand preparation can be relevant costs.

We propose studying XQDot4Z on Kuntur, a teaching-oriented RISC-V processor
with a five-stage pipeline. The aim is
to investigate a bounded operation and its costs, rather than to develop
a complete accelerator or execute a large language model at this stage.

\begin{figure}[t]
\centering
\begin{tikzpicture}[
 font=\scriptsize,
 box/.style={draw,minimum height=5.2mm,inner xsep=3pt,align=center},
 lbl/.style={font=\scriptsize\bfseries},
 arr/.style={-{Stealth[length=1.4mm]},shorten >=1pt,shorten <=1pt}]

\node[lbl] at (-3.55,0.55) {\ifSpanish Entrada\else Input\fi};
\node[box,text width=17mm] (in) at (-3.55,0)
 {\ifSpanish 8 pesos U4\\4 activaciones S8\else 8 U4 weights\\4 S8 activations\fi};

\foreach \y/\name/\code in {-1.25/b1/B1, -2.15/b2/B2, -3.05/b3/B3, -3.95/d/D}{
  \node[lbl] at (-3.55,\y) {\code};
}

\node[box,text width=20mm] (b1a) at (-1.35,-1.25)
 {\ifSpanish desplazar y sumar\else shift and add\fi};
\node[box,text width=20mm] (b1b) at (1.15,-1.25)
 {\ifSpanish restar 4 veces\else subtract 4 times\fi};

\node[box,text width=20mm] (b2a) at (-1.35,-2.15) {4 $\times$ MUL};
\node[box,text width=20mm] (b2b) at (1.15,-2.15)
 {\ifSpanish restar 4 veces\else subtract 4 times\fi};

\node[box,text width=20mm] (b3a) at (-1.35,-3.05) {XQDot4 $\rightarrow P$};
\node[box,text width=20mm] (b3b) at (1.15,-3.05) {$P - z\,S_a$};

\node[box,text width=20mm,very thick] (da) at (-1.35,-3.95) {XQDot4Zi};
\node[box,text width=20mm,draw=none] (db) at (1.15,-3.95)
 {\ifSpanish (ya corregido)\else (already corrected)\fi};

\node[box,text width=12mm] (out) at (3.35,-2.6) {$d$ (INT32)};

\foreach \a/\b in {b1a/b1b, b2a/b2b, b3a/b3b}{\draw[arr] (\a) -- (\b);}
\draw[arr] (da) -- (db);
\foreach \n in {b1a,b2a,b3a,da}{\draw[arr] (in.east) -| ($(\n.west)+(-3mm,0)$) -- (\n.west);}
\foreach \n in {b1b,b2b,b3b,db}{\draw[arr] (\n.east) -- ++(3mm,0) |- (out.west);}

\node[align=left,text width=26mm,anchor=west] at (2.1,-4.75)
 {\ifSpanish $S_a$ se calcula una vez por grupo y se reutiliza en las $N$ filas.
  \else $S_a$ is computed once per group and reused across the $N$ rows.\fi};
\end{tikzpicture}
\caption{\ifSpanish Trabajo que hace cada variante sobre un mismo grupo de cuatro
pesos. B1--B3 son puntos de comparación; D es el diseño en estudio. Las cuatro
leen los mismos datos empacados. La diferencia no está en el producto sino en
quién paga la corrección del punto cero: B1 y B2 la restan término a término,
B3 la factoriza y la amortiza entre filas, y D la absorbe dentro de la propia
instrucción.
\else Work each variant performs on one group of four weights. B1--B3 are points
of comparison; D is the design under study. All four read the same packed data.
The difference is not in the product but in who pays for the zero-point
correction: B1 and B2 subtract it term by term, B3 factors it and amortizes it
across rows, and D absorbs it inside the instruction itself.\fi}
\label{fig:comparison}
\end{figure}

\section{Background and Related Work}
In affine quantization, an integer code \(q\) approximately represents
\(r=s(q-z)\), where \(s>0\) is the scale and \(z\) the zero-point.
With activations whose zero-point is zero, the arithmetic kernel under
study is
\begin{equation}
d=\sum_{i=0}^{3}(q_{w,i}-z_w)q_{a,i}.
\label{eq:dot}
\end{equation}
Weights are stored as U4 (unsigned four-bit integers) and activations as
S8 (signed eight-bit integers). The proposed output is an INT32 subtotal
(a signed 32-bit integer); scales, bias, and accumulation of other
subtotals remain outside this operation.

XpulpNN and Dustin provide precedents for low-precision instructions and
mixed-precision arithmetic on RISC-V~\cite{xpulpnn,dustin}. Low-bit dot
products are therefore not treated as a new concept. The specific
contribution must be established through a more detailed comparison with
prior work and the project's own results.

An important alternative to fusion is
\begin{equation}
d=\sum_i q_{w,i}q_{a,i}-z_w\sum_i q_{a,i}.
\label{eq:factor}
\end{equation}
The activation sum can be reused across matrix rows. A valid comparison
must allow this optimization.

\section{Research Question and Scope}
\textit{Under what conditions does fusing zero-point correction into a
U4-by-S8 dot product reduce execution cost compared with scalar multiplication
and factored correction, accounting for activation reuse and zero-point
delivery?}

The initial scope is one core, batch size one, and small dot-product and
matrix--vector kernels. The first phase will evaluate correctness and
simulation cycles. Synthesis and board-level evaluation are later stages.
Energy or physical execution time will not be inferred solely from
instruction or cycle counts.

\section{Proposed Method}
\subsection{Platform and Verification}
The original processor is preserved, and development takes place in an
independent clone. Core corrections are verified separately and will form
a common baseline for the variants. Unimplemented instructions must be
detected, without claiming full conformance to an instruction set that
the processor does not yet cover.

The v0.1 numerical contract defines a 32-bit word \(W\) containing eight U4
weights and a word \(A\) containing four S8 activations. A selector \(h\)
chooses the lower or upper four weights of \(W\), without changing the
activations. Index zero occupies the least significant bits, and \(z_w\)
is shared by all four terms. The subtotal lies in \([-7680,7680]\) and is
represented in 32-bit two's complement, without implicit accumulation.

A Python reference using unbounded integers is available, checked through
examples, packing tests, exhaustive single-term tuples, and random dot
products compared with~\eqref{eq:factor}. The isolated RTL unit is combinational,
with four lanes, S13 products, and an adder tree with explicit sign extension.
It is verified in Icarus Verilog and Verilator. The experimental XQDot4Zi
interface optionally integrates it into Kuntur, using two source registers
and immediate \(z_w,h\). It reuses forwarding and carries metadata with the
instruction. It adds no dedicated wait in E, but physical timing has not been
checked. Runtime-loaded zero-points remain unresolved; an immediate cannot
replace them at no cost.

The baseline includes optional scalar MUL, independent of XQDot4Zi, returning
the low 32 bits of the product~\cite{riscvm}. It does not implement full M or
Zmmul. It is combinational, reuses forwarding, and is cancelled by flush/reset;
functional availability does not establish physical timing.

The packed comparator (B3) has a four-lane unit without zero-point
subtraction, using the same packing and selector as the fused XQDot4Z
variant (D).
Its subtotal lies in \([-7680,7620]\); the correction in~\eqref{eq:factor}
is outside this unit. The optional XQDot4 instruction integrates it using
two source GPRs, immediate \(h\), and the same forwarding. With weights set
to one, it also computes activation sums that can be reused across rows.

\subsection{Planned Comparison}
D will be compared with scalar MUL (B2) and, centrally, B3 with the correction
in~\eqref{eq:factor} (Fig.~\ref{fig:comparison}). B1 provides context without MUL:
it will use bounded-width masked bit decomposition, with its policy and code
frozen before the pilot. It does not represent optimal software. All variants
will share data and core corrections. Loads, loops, correction, and metadata
will be counted in addition to the arithmetic kernel.

The initial manifest sweeps row counts with one group of 32 elements and
separates constant controls from U4 codes balanced across phases. It distinguishes
zero, one, other powers of two, and remaining codes; mixed-program times are
not apportioned to these strata by frequency. It evaluates static specialization,
without extrapolation to dynamic metadata. Unrolling, register allocation,
scheduling, and correction reuse are fixed before measurement and shared by all
four variants. Under that schedule B3 corrects with one subtraction per row when
z repeats and with a multiply and a subtraction when it varies: the difference
comes from policy and schedule, not architecture. All cases,
including failures and slowdowns, will be reported. All four variants now have
verified kernels: row traversal is derived from each encoding, and only D is
forced to emit one row body per row when the zero point varies, which is
reported as a required floor rather than a choice.
Seeds cover inputs, not statistical repetitions; the pilot will examine
control, addresses, and stalls, not only equal cycle totals.

\balance
\section{Results and Discussion}
Isolated XQDot4Z passes 556\,734 comparisons per simulator, eight negative
controls, and four Icarus mutations; lint requires no waivers. Exhaustiveness
only covers single-term tuples. Immediate integration passes 25 programs,
including 1056 numerical operations, comparing retirement, writeback, stores,
and operands, and detecting two mutations. MUL passes 86\,704 vectors per
simulator and 49 programs with four configurations, eight negative controls
per simulator, and four Icarus mutations. Artifacts, logs, and hashes of these
test suites are retained.

The registered campaign ran \CampaignCases\ cases, all passing in both
simulators with none missing. Cycles are identical across the ten seeds in all
\CampaignCombos\ variant, shape and setting combinations, and across the sixteen
ZS phases, while the outputs differ in every one. That does not prove
independence from arbitrary data, nor does it reduce the grid: the reduction
gate is a deliberate revision and has not been taken.

Table~\ref{tab:ratio} holds the central comparison. D's advantage over B3 falls
as $N$ grows in every regime where B3 has correction to amortize, which is the
trend H2 anticipated. At $z{=}0$ there is no advantage at all: B3 omits Sa and
the correction, and the two variants coincide exactly. At $N{=}1$ a varying
zero point is degenerate, since a single row cannot have a $z$ that varies.

Under ZS the raw gap mixes two things. D is forced to emit one row body per row
because its $z$ is an immediate, while B3 loops, so part of the difference is
control rather than correction. Subtracting the structural twin, of the 164
cycles at $N{=}16$ some 47 belong to loop control and 117 to the correction:
the corrected ratio moves from 1.38 to 1.20 rather than 1.46 to 1.28.

Against scalar MUL, D uses between 6.2 and 6.8 times fewer cycles at $N{=}16$;
B1, without a multiplier, between 31 and 34 times more than D. That forced
specialization is not free: at $N{=}16$ with a varying $z$, D occupies 516 code
words against 72 for B3, because its encoding demands sixteen row bodies. The
cycle advantage and the text cost are the same interface decision.

These figures are cycles of a simulated in-order pipeline. They are not physical
time, energy or area, and D's additional hardware is not accounted for here:
that belongs to synthesis. The scope is static specialization; the regime with
zero points read at run time remains unresolved.

\begin{table}[t]
\caption{B3 cycles divided by D cycles, by regime and row count. Run
\CampaignRun, \CampaignCases\ cases; regimes are never pooled.}
\label{tab:ratio}
\centering
\begin{tabular}{crrrr}
\hline
$N$ & ZC $z{=}0$ & ZC $z{=}8$ & ZC $z{=}15$ & ZS \\
\hline
1 & 1.00 & 1.86 & 1.88 & 1.00 \\
4 & 1.00 & 1.25 & 1.26 & 1.46 \\
16 & 1.00 & 1.08 & 1.08 & 1.28 \\
\hline
\end{tabular}
\end{table}

The packed unit without correction passes 67\,910 vectors per simulator,
with ten negative controls per simulator and four Icarus mutations.
From 1680 RTL responses, 336 row/group outputs are reconstructed across
24 matrix--vector cases, reusing activation sums across rows.
In that isolated test suite, accumulation and correction run in the verifier.

XQDot4 integration passes 123 programs in both simulators, covering all eight
configurations of MUL and the two packed instructions, and detects three
Icarus mutations. Four paired two-row fixtures produce eight matching outputs
in B3 and D. In B3, Kuntur computes one activation sum, reuses it, and applies
MUL/SUB correction; D accumulates fused subtotals. These are correctness tests
with zero-points known at code-generation time, not optimized benchmarks or
dynamic-metadata tests.

\textit{Performance evaluation pending.} These results do not establish
acceleration, resource usage, physical timing, or efficient dynamic metadata delivery.

\subsection{Threats to Validity}
The shared skeleton reduces optimization bias but does not establish
optimality or eliminate software--architecture interaction.
The synthetic distribution does not necessarily represent a real network.
With z=0, separate correction can be omitted without establishing equal
performance. Dispatch among immediates would support dynamic z without
a new encoding, but still has no implementation or measured cost. Counters
require validation; RTL cycles do not establish physical timing.

\section{Conclusion}
Fusion reduces cycles against both factored correction and scalar
multiplication, but not uniformly: the benefit narrows as $N$ grows, vanishes at
$z{=}0$, and with a varying zero point is charged back in code text. The answer
to the question is not a number but a map of conditions, which is what this work
contributes.

Left outside are the regime with zero points read at run time, which this
interface cannot express, and physical evidence: area, frequency and energy need
synthesis and a different methodology. The fused variant's additional hardware
is not accounted for in these figures.
